% ══════════════════════════════════════════════════════════════
%  SECTION 2: INSTITUTIONAL BACKGROUND
% ══════════════════════════════════════════════════════════════

\section{Institutional Background}
\label{sec:background}


% ──────────────────────────────────────────────────────────────
\subsection{The Policy: An Enforcement Shift, Not a New Rule}
\label{sec:background:policy}

The rule at the center of this paper has existed, essentially unchanged,
since 1886. Law~7 of the Laws of the Game---the sole province of the
International Football Association Board (IFAB)---requires that referees
add time at the end of each half to compensate for time lost to
substitutions, assessment and removal of injured players, time-wasting,
disciplinary sanctions, stoppages for drinks or cooling breaks, VAR
interventions, and ``any other cause.'' The law does not cap the amount
of time that may be added. It does not instruct referees to exercise
discretion conservatively. It simply requires that lost time be
restored \citep{ifab2023laws}.

For decades, referees ignored this requirement in a remarkably uniform
way. The convention that emerged across European professional football
was to add between one and four minutes in the first half and between
three and five minutes in the second---regardless of how much time had
actually been consumed by interruptions. Empirical studies of ball-in-play
time consistently found that the ball was in active play for only 55--62
minutes per 90-minute match, implying 28--35 minutes of dead time per
match, of which referees restored fewer than eight
\citep{lames2024goal, wei2024goal}.
The gap between the rule on the books and the rule in practice was not
small and not contested; it was an equilibrium. Referees, federations,
broadcasters, and clubs operated within a shared understanding that
stoppage time was a compressed, approximate gesture rather than a literal
accounting of lost minutes.

What changed in November 2022 was not the law but its enforcement.
At the FIFA World Cup in Qatar, Pierluigi Collina, chairman of FIFA's
Referees Committee and the most influential figure in global match
officiating, instructed officials to add the \textit{actual} time lost
during each half. The results were immediately visible. England--Iran
ran to 117~total minutes of play. Several group-stage matches exceeded
100~minutes. The average second-half stoppage time at the tournament was
roughly double the pre-tournament convention across Europe's top leagues.
These matches were watched by cumulative audiences exceeding five billion
viewers, creating a public, salient, and unambiguous signal about the new
enforcement expectation.

The formal institutional step came at IFAB's Annual General Meeting in
March 2023, where the board endorsed the practice in characteristically
hortatory language: referees ``should'' add the actual time lost, not
that they ``must.'' This phrasing was deliberate. IFAB sets the Laws of
the Game but cannot compel national federations to adopt specific
enforcement standards. The directive thus left each federation with
discretion over implementation, creating the institutional heterogeneity
that drives our empirical design. Notably, UEFA explicitly opted out for
its club competitions: the Champions League and Europa League did not
adopt the new enforcement norm, providing a useful within-referee
placebo for future work.

The appropriate analogy is an enforcement crackdown, not a new statute.
The distinction matters for how we model the treatment. In the regulatory
compliance literature, \citet{greenstone2004did} showed that the Clean
Air Act's designation of ``nonattainment'' counties---which triggered
enforcement of pre-existing air quality standards---reduced pollution
substantially, even though the underlying legal standard had not changed.
\citet{johnson2020compliance} documented how shifts in enforcement
posture can alter behavior across an entire regulated population, even
absent changes in formal penalties. The FIFA directive fits squarely in
this tradition: the legal standard was constant, enforcement expectations
shifted, and the regulated population---professional referees---adjusted
behavior rapidly and almost universally.


% ──────────────────────────────────────────────────────────────
\subsection{The Dose Gradient}
\label{sec:background:dose}

A standard difference-in-differences design would partition our 15
leagues into treated and control groups and estimate the differential
change. We do this---designating the five Big~Five top-flight leagues
as treated and the remaining eleven as controls---and the estimate is
statistically significant. But the data make clear that this partition
obscures the directive's most important feature: it produced not a
binary treatment but a dose gradient, with enforcement intensity
proportional to institutional proximity to FIFA's refereeing apparatus.

Table~\ref{tab:enforcement} reports the change in mean second-half
stoppage time and the standardized effect size (Cohen's $d$) for each
league, sorted by the magnitude of the shift.

\begin{table}[t]
\centering
\caption{Enforcement Intensity by League}
\label{tab:enforcement}
\small
\begin{threeparttable}
\begin{tabular}{llccccl}
\toprule
League & Code & Pre Mean & Post Mean & $\Delta$ & Cohen's $d$ & Intensity \\
\midrule
Premier League$^\dagger$ & E0 & 5.39 & 7.50 & +2.11 & 1.13 & High \\
Ligue 1$^\dagger$ & F1 & 4.71 & 6.50 & +1.79 & 1.05 & High \\
Bundesliga$^\dagger$ & D1 & 4.34 & 6.12 & +1.78 & 0.92 & High \\
Scottish Premiership & SC0 & 2.83 & 4.30 & +1.47 & 0.83 & High \\
La Liga$^\dagger$ & SP1 & 5.18 & 6.64 & +1.46 & 0.70 & High \\
Belgian Pro League & B1 & 3.09 & 4.40 & +1.31 & 0.66 & High \\
2. Bundesliga & D2 & 3.95 & 4.98 & +1.04 & 0.59 & High \\
Championship & E1 & 5.67 & 6.70 & +1.03 & 0.59 & High \\
Primeira Liga & P1 & 5.97 & 6.94 & +0.97 & 0.42 & Moderate \\
Ligue 2 & F2 & 2.81 & 3.76 & +0.95 & 0.56 & Moderate \\
Greek Super League & G1 & 3.58 & 4.41 & +0.82 & 0.38 & Moderate \\
Serie A$^\dagger$ & I1 & 5.18 & 5.91 & +0.73 & 0.40 & Moderate \\
Eredivisie & N1 & 5.24 & 5.71 & +0.47 & 0.23 & Low \\
Serie B & I2 & 4.05 & 3.97 & -0.08 & -0.04 & Low \\
\bottomrule
\end{tabular}
\begin{tablenotes}\small
\item \textit{Notes:} Pre = seasons before 2023-24. Post = 2023-24 onward. $\Delta$ = Post $-$ Pre mean second-half stoppage time (minutes). Cohen's $d$ = standardized effect size. Intensity classified as High ($\Delta > 1.0$), Moderate ($0.5 < \Delta \leq 1.0$), Low ($\Delta \leq 0.5$). $^\dagger$ denotes Big Five league.
\end{tablenotes}
\end{threeparttable}
\end{table}


The pattern is striking.
The three largest responses came from Big~Five treated leagues: the
Premier League added 2.11~minutes ($d = 1.13$, $N_{\text{pre}} = 2{,}276$,
$N_{\text{post}} = 1{,}017$), Ligue~1 added 1.79~minutes ($d = 1.05$),
and the Bundesliga added 1.78~minutes ($d = 0.92$). These are large,
precisely estimated shifts in enforcement behavior, and they are
consistent with the narrative that the Big~Five leagues, which supply
the majority of referees to FIFA tournaments and maintain the closest
institutional ties to Collina's refereeing apparatus, responded most
strongly.

But the gradient does not respect the treated-control boundary. The
fourth-largest shift belongs to the Scottish Premiership---a control
league---which added 1.47~minutes ($d = 0.83$), outpacing La~Liga's
1.46~minutes ($d = 0.70$). The Belgian Pro League, another control,
shifted by 1.31~minutes ($d = 0.66$), exceeding Serie~A's 0.73~minutes
($d = 0.40$) by a wide margin. The 2.~Bundesliga ($+1.04$ minutes,
$d = 0.59$) and the English Championship ($+1.03$, $d = 0.59$)---the
paper's two most important within-country control leagues---shifted
substantially. The Primeira Liga ($+0.97$, $d = 0.42$), Ligue~2
($+0.95$, $d = 0.56$), and the Greek Super League ($+0.82$, $d = 0.38$)
all registered meaningful increases. The Eredivisie responded modestly
($+0.47$, $d = 0.23$). Only Serie~B showed no response at all
($-0.08$ minutes, $d = -0.04$). Thirteen of fourteen leagues shifted
in the same direction. This is not a pattern that a binary treatment
assignment can capture.

What predicts a league's position on this gradient? Four institutional
features emerge as explanatory. First, direct institutional ties to
FIFA's refereeing bodies. Leagues whose referees regularly officiate
FIFA-sanctioned tournaments had the strongest response, consistent with
those referees receiving the enforcement signal directly from Collina's
apparatus and carrying it back to domestic competition. Second, media
scrutiny and broadcasting pressure. The Premier League, Ligue~1, and the
Bundesliga are among the most-watched domestic competitions in the world;
deviation from the new norm would have been immediately visible and
widely discussed. The Scottish Premiership, while smaller, operates in
the same English-language media environment as the Premier League, which
may explain its unexpectedly strong response. Third, shared referee pools
within countries. Championship referees in England are drawn from a
pipeline that feeds into PGMOL, the body that manages Premier League
officiating; 2.~Bundesliga referees participate in a national pool
overseen by the DFB. To the extent that the new norm diffused through
professional networks rather than formal directives, within-country
second divisions would absorb the signal through referee pool overlap.
Fourth, professional refereeing culture. Leagues with full-time
professional referees---the Premier League has employed referees on
professional contracts since 2001---responded more strongly than those
relying on part-time officials, suggesting that the norm shift operated
through professional identity and peer expectations rather than formal
monitoring \citep[cf.][on referee responsiveness to institutional
incentives]{garicano2005favoritism}.

Taken together, the dose gradient implies that the directive functioned
as a global norm shift with heterogeneous uptake, not as a policy
applied to five leagues and withheld from eleven. The control leagues in
our DiD are not untreated; they are lower-dose. This interpretation has
direct consequences for the empirical strategy we develop in
Section~\ref{sec:strategy}: a binary DiD estimates the differential
response between high- and medium-dose leagues, yielding a lower bound
on the directive's total effect. The summary statistics reinforce the
point. In the pre-directive period, treated leagues averaged 4.99~minutes
of second-half stoppage time ($N = 10{,}763$) and control leagues
averaged 4.32~minutes ($N = 15{,}865$). In the post-directive period,
treated leagues rose to 6.56~minutes ($+1.57$) and control leagues rose
to 5.62~minutes ($+1.30$). A control-group increase of 1.30~minutes---82\%
as large as the treated-group increase---confirms that the norm diffused
well beyond the five designated treatment leagues.


% ──────────────────────────────────────────────────────────────
\subsection{The Serie~A Problem}
\label{sec:background:seriea}

Serie~A occupies an uncomfortable position in the empirical design. It
is classified as treated on defensible institutional grounds: it belongs
to the Big~Five, the Italian football federation (FIGC) issued explicit
guidance to referees about the directive, and Italian referees
officiate regularly at FIFA tournaments. Yet Serie~A's response was
muted. The league added only 0.73~minutes of second-half stoppage time
($d = 0.40$), placing it twelfth of fourteen leagues---closer to the
Greek Super League and the Eredivisie than to the Premier League or
Ligue~1.

Match-level compliance statistics sharpen the contrast. In the
post-directive period, only 56.4\% of Serie~A matches exceeded the
league's own pre-directive median stoppage time, compared with 81.9\%
for the Premier League. Where the English game exhibited a sharp,
near-universal upward shift in the distribution of stoppage time,
Serie~A's distribution shifted modestly and with considerable overlap
between the pre- and post-directive periods.

The muted response has consequences for identification. In the
Goodman-Bacon decomposition of our binary DiD, the bilateral
estimate for the Italy pair (Serie~A vs.\ Serie~B) is \textit{negative},
because Serie~B---the paired control---showed no response at all.
The Serie~A--Serie~B comparison thus contributes negatively to the
aggregate treatment effect, mechanically attenuating the overall
estimate. This is not a failure of the decomposition; it is an honest
reflection of the fact that the within-Italy contrast does not support
the hypothesis that treated leagues uniformly outpaced their
within-country controls.

We adopt an honest approach: all main specifications are reported both
with and without Serie~A, and the dose-response framework developed in
Section~\ref{sec:strategy} allows Serie~A to occupy its natural position
on the enforcement gradient rather than forcing it into a binary
category. Several explanations for the muted response are plausible---a
more conservative Italian refereeing culture, a different institutional
relationship between the FIGC and its match officials, or less intense
media pressure for compliance---but we do not attempt to adjudicate
among them. What matters for identification is that the anomaly is
transparent and that the empirical strategy accommodates it.


% ──────────────────────────────────────────────────────────────
\subsection{Within-Country Institutional Structure}
\label{sec:background:within}

The strongest element of our research design is the availability of
within-country league pairs: top-flight treated leagues paired with
second-division control leagues that share a national federation,
regulatory environment, and (to varying degrees) a referee pool, but
differ in their institutional proximity to FIFA and in the intensity of
media and broadcasting scrutiny. These pairs difference out
country-level confounders---national footballing culture, federation
organizational capacity, macroeconomic conditions---leaving variation
in enforcement intensity driven by the within-country institutional
differences that the dose-gradient logic predicts.

England provides the cleanest pair. The Premier League and the
Championship share the Football Association as their governing body but
operate under different refereeing organizations. The Premier League's
referees are managed by the Professional Game Match Officials Limited
(PGMOL) and have been full-time professionals since 2001. Championship
referees are managed by the EFL and are not full-time professionals,
though the most promising are promoted to the PGMOL pathway. The
Premier League's referees thus had direct exposure to the FIFA signal
through international assignments and PGMOL's institutional relationship
with FIFA; Championship referees received the signal indirectly, filtered
through professional networks and the broader media environment.
The difference in response---$+2.11$ versus $+1.03$ minutes---is
consistent with this institutional channel, though the Championship's
substantial shift confirms that the norm penetrated well beyond the
directly treated tier.

Germany offers a similar structure. The Bundesliga and 2.~Bundesliga
share DFB governance but maintain different refereeing pools. The
Bundesliga's shift ($+1.78$ minutes) exceeded the 2.~Bundesliga's
($+1.04$), but the second division's response was again large enough
to constitute meaningful contamination of the control group. France
presents a pair with more institutional overlap: Ligue~1 and Ligue~2
share some referees, as the Direction Technique de l'Arbitrage manages
both competitions. Ligue~1 shifted by 1.79~minutes and Ligue~2 by
0.95---a gap that may partially reflect the direct referee-pool
overlap rather than institutional distance alone.

Italy, as discussed above, is the anomalous pair. Serie~A and Serie~B
operated under different referee assignment systems: the Commissione
Arbitri Nazionali (CAN) managed Serie~A assignments, while Serie~B
used a separate assignment body. This institutional separation may help
explain why the directive's signal, already weakly received in Serie~A,
failed to reach Serie~B entirely. Spain is absent from the within-country
comparison because Segunda Divisi\'{o}n data are unavailable for the
post-2022--23 seasons, precluding a parallel analysis.

The within-country pairs serve a dual purpose. In the binary DiD, they
provide four tightly matched comparisons that control for national-level
confounders. In the dose-response framework, they illuminate the
mechanism of norm transmission: the enforcement signal attenuated
predictably as it passed from FIFA through national federations to
top-flight leagues and, more weakly, to second divisions. The
gradient within countries mirrors the gradient across countries,
reinforcing the interpretation that institutional proximity to the
directive's source---not formal legal designation---determined
enforcement intensity.
